% Appendix: entanglement proof details.

\section{Entanglement proof details}\label{app:entangle-proof}

This appendix records the technical details of the proof of
Proposition~\ref{prop:entanglement} that were compressed in the main body, and
notes the exact role of the continuity hypothesis on the MPNN's
update and readout.

\subsection{The colour trace on a monochromatic triangle}

Consider a sample $(\gph_A, x_A) \in \mathrm{supp}(\mathcal{M}_n^A)$
and a vertex $v$ in a $(+1)$-triangle $T = \{v, u_1, u_2\}$. The
initial colour of $v$ is $c_v^{(0)} = \tau(+1)$; call this $\alpha_+$.
By the definition of the ensemble, $x_{u_1} = x_{u_2} = +1$ as well,
so $c_{u_1}^{(0)} = c_{u_2}^{(0)} = \alpha_+$.

At round $1$,
\[
c_v^{(1)} = \tau\!\left(c_v^{(0)},\ \{\!\{c_{u_1}^{(0)}, c_{u_2}^{(0)}\}\!\}\right) = \tau(\alpha_+, \{\!\{\alpha_+, \alpha_+\}\!\}).
\]
Denote this colour $\beta_+$. By the same computation, $c_{u_1}^{(1)} =
c_{u_2}^{(1)} = \beta_+$.

At round $2$, every neighbour of $v$ has colour $\beta_+$, so
\[
c_v^{(2)} = \tau(\beta_+, \{\!\{\beta_+, \beta_+\}\!\}) = \gamma_+,
\]
and the same holds for $u_1$ and $u_2$. From round $2$ onward, every
vertex of the triangle shares a common colour that evolves as
$\gamma_+ \mapsto \tau(\gamma_+, \{\!\{\gamma_+, \gamma_+\}\!\})$; the
evolution is one-to-one, so no two rounds produce the same colour on
the triangle, but the partition of vertices induced by refinement is
the same partition (the singleton class containing all three vertices
of the triangle). Refinement has stabilised at round $1$ in the sense
that the partition does not refine further.

The same computation applies to every $(+1)$-triangle in the sample,
so every vertex in every $(+1)$-triangle receives the same stable
colour. Analogously, every vertex in every $(-1)$-triangle receives a
common colour $\beta_-$. There are exactly $n/2$ vertices of each
kind.

\subsection{The colour trace on a monochromatic six-cycle}

Consider a sample $(\gph_B, x_B) \in \mathrm{supp}(\mathcal{M}_n^B)$
and a vertex $v$ in a $(+1)$-cycle $C = \{v_0, v_1, v_2, v_3, v_4,
v_5\}$ with edges $\{v_i, v_{i+1 \bmod 6}\}$. Each $v_i$ has
$x_{v_i} = +1$ and two neighbours $v_{i \pm 1 \bmod 6}$, both with
feature $+1$.

At round $1$, $c_{v_i}^{(1)} = \tau(\alpha_+, \{\!\{\alpha_+,
\alpha_+\}\!\}) = \beta_+$: the same colour that a vertex of a
$(+1)$-triangle receives at round $1$.

At round $2$, every neighbour of $v_i$ has colour $\beta_+$, so
$c_{v_i}^{(2)} = \tau(\beta_+, \{\!\{\beta_+, \beta_+\}\!\}) =
\gamma_+$: the same colour that a vertex of a $(+1)$-triangle
receives at round $2$.

Higher rounds continue in the same one-to-one evolution: the partition
does not refine further. The graph-level fingerprint agrees with the
fingerprint on samples from $\mathcal{M}_n^A$:
\[
\mathrm{hist}(\gph_A, x_A) = \mathrm{hist}(\gph_B, x_B) = \{\!\{\gamma_+^{n/2},\ \gamma_-^{n/2}\}\!\}.
\]

\subsection{Continuity of the MPNN's factorisation}

The proof of Proposition~\ref{prop:entanglement}(i) invokes Proposition~\ref{prop:upper} to
pass from $\mathrm{hist}(\gph_A, x_A) = \mathrm{hist}(\gph_B, x_B)$
to $M(\gph_A, x_A) = M(\gph_B, x_B)$ for every $M \in \mpnn$ with
continuous update and readout.

Proposition~\ref{prop:upper} in fact shows this stronger statement: $M(\gph_A,
x_A)$ is a function of $\mathrm{hist}(\gph_A, x_A)$ alone.
Continuity of the update and readout does not enter into whether
this function is well-defined; it enters only when one considers
uniform approximation, which the ensemble-level theorem does not
require. On samples with matching fingerprint, the MPNN's output is
pointwise equal regardless of continuity assumptions.

\subsection{Extension to noisy features}

The construction of Proposition~\ref{prop:entanglement} assigns features by
component: every vertex in a $(+1)$-triangle receives $x_v = +1$.
This is the exact-alignment case.

If one applies a feature-flip perturbation with probability $\sigma$
to each vertex independently, the two ensembles $\mathcal{M}_n^A(\sigma)$
and $\mathcal{M}_n^B(\sigma)$ are no longer $1$-WL equivalent
pointwise: a $(+1)$-triangle after perturbation may contain a vertex
with $x_v = -1$, and the round-one $1$-WL colour of such a vertex
depends on the number of $+1$ vs. $-1$ neighbours in the triangle,
which differs from the corresponding count in a $6$-cycle whose
vertices have been perturbed similarly.

Numerical simulation of the two ensembles at $\sigma = 0.05$ and
$n = 300$ gives a depth-$1$ $1$-WL total variation of approximately
$0.14$ between the two ensembles. The impossibility as stated
requires $\sigma = 0$.

\subsection{Extension to weighted rewiring within components}

The ensembles $\mathcal{M}_n^A$ and $\mathcal{M}_n^B$ use uniform
random assignment of vertices to triangles or $6$-cycles. Any
alternative assignment procedure (weighted random, deterministic,
or dependent on an auxiliary feature) preserves the theorem, as
long as the resulting graph decomposes into monochromatic triangles
or monochromatic $6$-cycles: the proof depends only on the
component structure of the featured graph.

Extensions that connect the components (adding cross-component
edges) break the theorem. If cross-triangle edges are added to
samples from $\mathcal{M}_n^A$, the round-one $1$-WL colour of a
vertex touched by a cross-triangle edge sees a neighbour outside
its triangle. The neighbour's feature may or may not match the
vertex's, but in either case the colour differs from a vertex in a
purely monochromatic triangle. Cross-cycle edges added to samples
from $\mathcal{M}_n^B$ have the same effect. Matching the two
ensembles with cross-component edges requires the additional edges
to have matched sign-block statistics, and this is the SBM
extension analysed in Section~\ref{sec:obstacle}.
